Moreover, the preceding decomposition induces isomorphisms
\[
\underline{\operatorname{Dyn}}_t(G)\simeq\underline{\operatorname{Dyn}}(G_t),\qquad
\underline{\Aut}_{S\text{-gr.}}(G)\simeq\prod_{t\in\mathbf T}\underline{\Aut}_{S\text{-gr.}}(G_t).
\]
\end{proposition}
\begin{proof}
This has indeed been proved above when $G$ is split. In the general case, one may suppose $G$ of constant type $\mathscr R$. Using the preceding decomposition of $A_S(\mathscr R)$ and 1.17, one deduces the desired decomposition of $G$. The other results are then proved by descent.
\end{proof}
\begin{remark}{5.6}\label{XXIV.5.6}
More generally, if $G$ and $H$ are two adjoint semisimple groups, one has canonical isomorphisms as follows (the diagram is commutative):
\[
\xymatrix@C=1.2em{
\underline{\Isom}_{S\text{-gr.}}(G,H)\ar[r]^{\sim}\ar[d]&\prod_{t\in\mathbf T}\underline{\Isom}_{S\text{-gr.}}(G_t,H_t)\ar[d]\\
\underline{\operatorname{Isomext}}(G,H)\ar[r]^{\sim}\ar[d]_{\wr}&\prod_{t\in\mathbf T}\underline{\operatorname{Isomext}}(G_t,H_t)\ar[d]^{\wr}\\
\underline{\Isom}_{\mathrm{Dyn}}(\underline{\operatorname{Dyn}}(G),\underline{\operatorname{Dyn}}(H))\ar[r]^{\sim}&\prod_{t\in\mathbf T}\underline{\Isom}_{\mathrm{Dyn}}(\underline{\operatorname{Dyn}}(G_t),\underline{\operatorname{Dyn}}(H_t)).}
\]
\end{remark}
\begin{remark}{5.7}\label{XXIV.5.7}
One can give an intrinsic characterization of $G_t$, which we state below without proof: it is the largest normal (and indeed characteristic) reductive subgroup of $G$ isotypic of type $t$.
\end{remark}
The preceding proposition allows one to reduce the study of adjoint semisimple groups to that of isotypic adjoint semisimple groups. It is this study that we shall carry out below.
\numpara{5.8}\label{XXIV.5.8}
Let $\mathscr R$ be a \emph{simple} reduced pinned adjoint root datum of type $t$, $I$ a nonempty finite set, and $\mathscr R^I$ the product root datum of copies $\mathscr R_i$ of $\mathscr R$, for $i\in I$. Denote by $E$ the group of automorphisms of $\mathscr R$, which identifies with the group of automorphisms of the Dynkin diagram $D$ of $\mathscr R$. The Dynkin diagram of $\mathscr R^I$ identifies with $D^I$, which shows that one has an exact sequence:
\[
1\to\Aut(D)^I\to\Aut(D^I)\to\Aut(I)\to1,
\]
where $\Aut(I)$ denotes the group of permutations of $I$. It follows that the canonical isomorphism
\[
\operatorname{\acute Ep}_S(\mathscr R)^I\simeq\operatorname{\acute Ep}_S(\mathscr R^I)
\]
induces an exact sequence
\[
1\to A_S(\mathscr R)^I\to A_S(\mathscr R^I)\to\underline{\Aut}(I_S)\to1,
\]
the last group being the constant $S$-group associated to $\Aut(I)$. Observe moreover that $I$ identifies canonically with the set of connected components of $D^I$.
If $G$ is a semisimple $S$-group of type $\mathscr R^I$, defining (cf. 1.17) a principal homogeneous bundle $P$ under $A_S(\mathscr R^I)$, the definition of $\underline{\operatorname{Dyn}}(G)$ by descent (3.7), and that of $\underline{\operatorname{Dyn}}_0(G)$ (5.2), show that the bundle associated to $P$ by the morphism $A_S(\mathscr R^I)\to\underline{\Aut}(I_S)$ is nothing other than $\underline{\Isom}_S(I_S,\underline{\operatorname{Dyn}}_0(G))$, corresponding to the form $\underline{\operatorname{Dyn}}_0(G)$ of $I_S$. Using again the equivalence of categories 1.17 and the preceding exact sequence, one deduces by a formal argument that there exist a reductive $\underline{\operatorname{Dyn}}_0(G)$-group of type $\mathscr R$, say $G_0$, and an $S$-isomorphism $G\simeq\prod_{\underline{\operatorname{Dyn}}_0(G)/S}G_0$.
\begin{proposition}{5.9}\label{XXIV.5.9}
Let $S$ be a scheme and $G$ an isotypic adjoint semisimple $S$-group of type $t$. There exist a simple adjoint semisimple $\underline{\operatorname{Dyn}}_0(G)$-group $G_0$ of type $t$, and an $S$-isomorphism (unique)
\[
G\simeq\prod_{\underline{\operatorname{Dyn}}_0(G)/S}G_0.
\]
Moreover, this isomorphism induces an exact sequence
\[
1\to\prod_{\underline{\operatorname{Dyn}}_0(G)/S}\underline{\Aut}_{S\text{-gr.}}(G_0)
\to\underline{\Aut}_{S\text{-gr.}}(G)\to\underline{\Aut}_S(\underline{\operatorname{Dyn}}_0(G))\to1.
\]
% typo?: source subscript S-gr. on Aut(G_0), though G_0 is over Dyn_0(G); kept as printed.
\end{proposition}
\begin{proof}
One may evidently suppose $G$ of constant type $nt$. The first assertion has already been proved (the uniqueness assertion is evident). The second then follows from the split case by descent.
\end{proof}
One can combine 5.5 and 5.9:
\begin{proposition}{5.10}\label{XXIV.5.10}
Let $S$ be a scheme, $G$ an adjoint semisimple $S$-group, and $\mathscr D=\underline{\operatorname{Dyn}}(G)$ its Dynkin scheme.
\begin{enumeratei}
\item There exists a canonical decomposition
\[
G\simeq\prod_{t\in\mathbf T}\prod_{\mathscr D_{0,t}/S}G_{0,t}\simeq\prod_{\underline{\operatorname{Dyn}}_0(G)/S}G_0,
\]
where each $G_{0,t}$ is a simple adjoint $\underline{\operatorname{Dyn}}_{0,t}(G)$-group of type $t$ (resp. where $G_0$ is an adjoint semisimple $\underline{\operatorname{Dyn}}_0(G)$-group whose type at $x\in\underline{\operatorname{Dyn}}_0(G)$ is $a(x)\in\mathbf T$ (the morphism $a:\underline{\operatorname{Dyn}}_0(G)\to\mathbf T_S$ was defined in 5.2)).
\item The preceding decompositions induce isomorphic exact sequences (written vertically), where $\underline{\Aut}_{S,a}(\mathscr D_0)$ denotes the scheme of automorphisms of $\underline{\operatorname{Dyn}}_0(G)$ which commute with the morphism $a:\underline{\operatorname{Dyn}}_0(G)\to\mathbf T_S$:
\[
\xymatrix@C=1.8em{
1\ar[d]&1\ar[d]\\
\prod_{t\in\mathbf T}\prod_{\mathscr D_{0,t}/S}\underline{\Aut}_{\mathrm{gr.}}(G_{0,t})\ar[r]^{\sim}\ar[d]&\prod_{\underline{\operatorname{Dyn}}_0(G)/S}\underline{\Aut}_{\mathrm{gr.}}(G_0)\ar[d]\\
\underline{\Aut}_{S\text{-gr.}}(G)\ar[r]\ar[d]&\underline{\Aut}_{S\text{-gr.}}(G)\ar[d]\\
\prod_{t\in\mathbf T}\underline{\Aut}_S(\underline{\operatorname{Dyn}}_{0,t}(G))\ar[r]^{\sim}\ar[d]&\underline{\Aut}_{S,a}(\underline{\operatorname{Dyn}}_0(G))\ar[d]\\
1&1}
\]
\end{enumeratei}
\end{proposition}
\begin{corollary}{5.11}\label{XXIV.5.11}
Under the preceding conditions, the following three categories are equivalent:
\begin{enumeratei}
\item the category of principal homogeneous bundles under $G$.
\item the category of principal homogeneous bundles under $G_0$.
\item the product category, for $t\in\mathbf T$, of the categories of principal homogeneous bundles under the $G_{0,t}$.
\end{enumeratei}
\end{corollary}
This follows formally from the preceding decompositions and from 8.4.
\begin{corollary}{5.12}\label{XXIV.5.12}
One has canonical isomorphisms
\[
\underline{\operatorname{Tor}}(G)\simeq\prod_{t\in\mathbf T}\prod_{\mathscr D_{0,t}/S}\underline{\operatorname{Tor}}(G_{0,t})\simeq\prod_{\underline{\operatorname{Dyn}}_0(G)/S}\underline{\operatorname{Tor}}(G_0),
\]
and likewise on replacing $\underline{\operatorname{Tor}}$ by $\underline{\operatorname{Bor}}$ (resp. $\underline{\operatorname{Kil}}$).
\end{corollary}
Trivial from the split case.
\begin{remark}{5.13}\label{XXIV.5.13}
The canonical morphism
\[
\underline{\operatorname{Dyn}}(G)\to\underline{\operatorname{Dyn}}_0(G)
\]
allows one to consider $\underline{\operatorname{Dyn}}(G)$ as a Dynkin $\underline{\operatorname{Dyn}}_0(G)$-scheme; in fact it is the Dynkin scheme $\underline{\operatorname{Dyn}}(G_0)$ of $G_0$.

Likewise, if $T\subset B$ is a Killing pair of $G$, canonically corresponding to the Killing pair $T_0\subset B_0$ of $G_0$, one verifies that the obstructions to quasi-splitting $G$ and $G_0$, which lie (3.9) in $\Pic(\underline{\operatorname{Dyn}}(G))=\Pic(\underline{\operatorname{Dyn}}(G_0))$, coincide. One deduces:
\end{remark}
\begin{corollary}{5.14}\label{XXIV.5.14}
The following conditions are equivalent:
\begin{enumeratei}
\item $G$ is quasi-splittable,
\item $G_0$ is quasi-splittable,
\item each $G_{0,t}$, $t\in\mathbf T$, is quasi-splittable.
\end{enumeratei}
\end{corollary}

\sgasection{6}{Automorphisms of Borel subgroups of reductive groups}
\begin{lemma}{6.1}\label{XXIV.6.1}
Let $S$ be a scheme, $(G,T,M,R)$ a split $S$-group, $\Delta$ a system of simple roots of $R$, and $B$ the corresponding Borel subgroup. Then $B^{\mathrm u}$ is generated as an (fppf) sheaf of groups by the $U_\alpha$, $\alpha\in\Delta$.
\end{lemma}
Let $H$ be the subsheaf of groups of $B^{\mathrm u}$ generated by the $U_\alpha$, $\alpha\in\Delta$. Let us prove $H\supset U_\beta$ ($\beta\in R_+$) by induction on the integer $\operatorname{ord}(\beta)=\operatorname{ord}_\Delta(\beta)>0$ (cf. Exp.~XXI, 3.2.15). The assertion holds by hypothesis for $\operatorname{ord}(\beta)=1$. Suppose therefore $\operatorname{ord}(\beta)>1$ and $U_\gamma\subset H$ whenever $\operatorname{ord}(\gamma)<\operatorname{ord}(\beta)$. There exists $\alpha\in\Delta$ such that $\beta-\alpha\in R_+$ (Exp.~XXI, 3.1.1). Let $p$ be the largest integer such that $\beta-p\alpha=\beta'\in R_+$. One has $U_\alpha,U_{\beta'}\subset H$, $\beta'-\alpha\notin R$. One is therefore reduced to proving:
\begin{lemma}{6.2}\label{XXIV.6.2}
Take up the notation of Exp.~XXIII, 6.4. Suppose $p=1$, i.e. $\beta-\alpha$ is not a root. If $H$ is a subsheaf of groups of $B^{\mathrm u}$ such that $U_\alpha,U_\beta\subset H$, then $U_{i\alpha+j\beta}\subset H$ whenever $i\alpha+j\beta\in R$, $i>0$, $j>0$.
\end{lemma}
\begin{proof}
Let us distinguish four cases according to the value of $q=0,1,2,3$. In what follows $x$ and $y$ denote two arbitrary sections of $\mathbf G_{\mathrm a,S'}$, $S'\to S$.

If $q=0$, there is nothing to prove.

If $q=1$, one has $p_{\alpha+\beta}(\pm x)=p_\beta(-y)p_\alpha(-x)p_\beta(y)p_\alpha(x)\in H(S')$, hence $U_{\alpha+\beta}\subset H$.
% typo?: the source prints +/- x, not +/- xy, for arbitrary x,y; kept as printed.

If $q=2$, one similarly has
\[
p_{\alpha+\beta}(\pm xy)p_{2\alpha+\beta}(\pm x^2y)=p_\beta(-y)p_\alpha(-x)p_\beta(y)p_\alpha(x)\in H(S'),
\]
whence, after changing certain signs,
\[
F(x,y)=p_{\alpha+\beta}(xy)p_{2\alpha+\beta}(x^2y)\in H(S').
\]
Since $p_{\alpha+\beta}$ and $p_{2\alpha+\beta}$ commute, one then has $F(x,1)F(1,-x)=p_{2\alpha+\beta}(x^2-x)$.\nde{We have corrected $F(-1,x)$ to $F(1,-x)$.}If $a\in\Ga(S)$, the equation $X^2-X=a$ defines a free surjective extension of $S$ (it is $\Spec\OS[X]/(X^2-X-a)$); one therefore has $p_{2\alpha+\beta}(a)\in H(S')$, hence $U_{2\alpha+\beta}\subset H$, hence also $U_{\alpha+\beta}\subset H$ (since $F(1,a)\in H(S')$).

If $q=3$, one similarly has
\[
F(x,y)=p_{\alpha+\beta}(xy)p_{2\alpha+\beta}(x^2y)p_{3\alpha+\beta}(x^3y)p_{3\alpha+2\beta}(x^3y^2)\in H(S');
\]
and since
\[
p_{3\alpha+2\beta}(\pm x)=F(1,1)^{-1}p_\beta(-x)F(1,1)p_\beta(x)\in H(S'),
\]
one obtains that $U_{3\alpha+2\beta}\subset H$ and therefore
\[
K(x,y)=p_{\alpha+\beta}(xy)p_{2\alpha+\beta}(x^2y)p_{3\alpha+\beta}(x^3y)\in H(S').
\]
Then computing
\[
K(x+y,1)K(-x,1)K(1,y)^{-1}
\]
modulo $U_{3\alpha+2\beta}$, one finds
\[
p_{2\alpha+\beta}(2x^2+y^2+2xy-y)p_{3\alpha+\beta}(y^3+3xy^2+3x^2y-y)\in H(S').
\]
If $a\in\Ga(S)$, the ``equations''
\[
\begin{aligned}
x^2&=-xy-y+1-a\\
y^2&=3y-2+3a
\end{aligned}
\]
define a free surjective extension of $S$; one then has $y^3+3xy^2+3x^2y-y=0$ and the preceding expression gives $p_{2\alpha+\beta}(a)\in H(S')$.\nde{We have corrected $p_{3\alpha+\beta}(a)$ to $p_{2\alpha+\beta}(a)$, and given the end of the argument in detail.}One has therefore proved that $H$ contains $U_{2\alpha+\beta}$ and $U_{3\alpha+2\beta}$ and that
\[
E(x,y)=p_{\alpha+\beta}(xy)p_{3\alpha+\beta}(x^3y)\in H(S').
\]
Since $p_{\alpha+\beta}(xy)$ and $p_{3\alpha+\beta}$ commute, one is reduced to the preceding computation, i.e. $E(1,x)E(1,-x)=p_{3\alpha+\beta}(x^3-x)$, whence $U_{3\alpha+\beta}\subset H$, then $U_{\alpha+\beta}\subset H$.
% typo?: E(1,x)E(1,-x) in the source; kept as printed.
\end{proof}
\begin{remark}{6.2.1}\label{XXIV.6.2.1}
The preceding proof shows that one could have replaced the hypothesis ``$H$ contains $U_\alpha$ and $U_\beta$'' by ``$H$ contains $U_\alpha$ or $U_\beta$, and $H$ is invariant under $\mathrm{int}(U_\alpha)$ and $\mathrm{int}(U_\beta)$''.
\end{remark}
\begin{theorem}{6.3}\label{XXIV.6.3}
Let $S$ be a scheme, $G$ and $G'$ two semisimple $S$-groups, and $B$ (resp. $B'$) a Borel subgroup of $G$ (resp. $G'$). Every isomorphism $B\isomto B'$ is induced by a unique isomorphism $G\isomto G'$.
\end{theorem}
\begin{proof}
The assertion is local for the étale topology and one may suppose $G$ split: $G=(G,T,M,R)$, $B$ being defined by the system $R_+$ of positive roots of $R$. The given isomorphism $u:B\isomto B'$ induces an isomorphism of $T$ onto a maximal torus $T'$ of $B'$, hence of $G'$. The given isomorphism $T\simeq\mathrm D_S(M)$ gives an isomorphism $T'\simeq\mathrm D_S(M)$ in which the elements of $R_+$ become the constant roots of $B'$ relative to $T'$, hence the elements of $R$ the constant roots of $G'$ relative to $T'$. Since $G$ and $G'$ are semisimple, the coroots are determined by duality (Exp.~XXI, 1.2.5), which proves that $(T',M,R)$ is a splitting of $G'$ such that $\mathscr R(G)=\mathscr R(G')$.

Applying Exp.~XXIII, 5.1 (Uniqueness theorem), one deduces that there exists a unique isomorphism $G\isomto G'$ agreeing with $u$ on $T$ and the $U_\alpha$, $\alpha\in\Delta$. By 5.1, the restriction of this isomorphism to $B$ is $u$.
% typo?: source reference 5.1 here; kept as printed.
\end{proof}
\begin{remark}{6.3.1}\label{XXIV.6.3.1}
(i) Using Exp.~XXII, 4.1.9 and reasoning as in loc. cit. 4.2.12, one can replace ``isomorphism'' in the statement of the theorem by ``isogeny'' (resp. ``central isogeny'').

(ii) The theorem is false for reductive groups. Take for example $G=G'=\SL_2\times\Gm$, identified with the group of the following matrices:
\[
\left\{\begin{pmatrix}a&b&0\\c&d&0\\0&0&h\end{pmatrix}\ \middle|\ ad-bc=1,\ h\text{ invertible}\right\};
\]
take for $B=B'$ the Borel subgroup defined by $c=0$ and for $u$ the following automorphism of $B$: $u(a,b,d,h)=(a,b,d,ah)$.
\end{remark}
\begin{corollary}{6.4}\label{XXIV.6.4}
The functor $(G,B)\mto B$ from the category consisting of pairs $(G,B)$, where $G$ is a semisimple $S$-group and $B$ a Borel subgroup of $G$, to the category of $S$-group schemes (the morphisms are the isomorphisms) is fully faithful.
\end{corollary}
\begin{corollary}{6.5}\label{XXIV.6.5}
Let $S$ be a scheme, $G$ a semisimple $S$-group, $B$ a Borel subgroup of $G$, and $B'$ an $S$-group locally isomorphic to $B$ for (fpqc). Then $B'$ is locally isomorphic to $B$ for the local finite étale topology\nde{Denoted by (étf) in Exp.~IV 6.3. In other words, for every $s\in S$ there exist an open neighborhood $U$ of $s$ and a finite surjective étale morphism $V\to U$ such that $B'\times_S V\simeq B\times_S V$.} and there exists a semisimple $S$-group $G'$ of which $B'$ is a Borel subgroup; $G'$ is unique up to a unique isomorphism inducing the identity on $B'$.
\end{corollary}
\begin{corollary}{6.6}\label{XXIV.6.6}
Let $S$ be a scheme, $G$ a semisimple $S$-group, and $B$ a Borel subgroup of $G$. Then $\underline{\Aut}_{S\text{-gr.}}(B)$ is representable by a smooth affine $S$-scheme, $\underline{\operatorname{Autext}}(B)$ is representable by a finite étale $S$-scheme, and the evident morphisms induce an isomorphism of exact sequences
\[
\xymatrix@C=1.8em{
1\ar[r]&B^{\mathrm{ad}}\ar[r]\ar[d]_{\wr}&\underline{\Aut}_{S\text{-gr.}}(G,B)\ar[r]\ar[d]_{\wr}&\underline{\operatorname{Autext}}(G)\ar[r]\ar[d]_{\wr}&1\\
1\ar[r]&B^{\mathrm{ad}}\ar[r]&\underline{\Aut}_{S\text{-gr.}}(G)\ar[r]&\underline{\operatorname{Autext}}(B)\ar[r]&1.}
\]
% typo?: bottom row has Aut(G) in the source, not Aut(B); kept as printed.
\end{corollary}
This follows at once from 2.1 and the preceding results. We leave to the reader the task of developing the analogues of the results of nos. 1, 2, 3, 4 in the above setting.
\begin{remark}{6.7}\label{XXIV.6.7}
If $S$ is a scheme and $B$ an $S$-group, $B$ can therefore be a Borel subgroup of only one semisimple group, completely determined by $B$. It therefore remains to characterize the $S$-groups $B$ which are indeed Borel subgroups of semisimple groups.\nde{One may ask whether every smooth affine $S$-group, each of whose geometric fibers is a Borel subgroup of a semisimple group, is a Borel subgroup of a semisimple $S$-group. (We have removed the ``counterexample'' given in the original, for $S=$ the scheme of dual numbers over a field $k$, which was erroneous, as M.~Demazure pointed out to us.)}
\end{remark}

\sgasection{7}{Representability of the functors $\underline{\Hom}_{S\text{-gr.}}(G,H)$, for reductive $G$}
\numpara{7.1}\label{XXIV.7.1}\textbf{The split case}
\numpara{7.1.1}\label{XXIV.7.1.1}
Let $S$ be a scheme, $G$ a pinned $S$-group, $T$ its maximal torus, $\Delta$ the set of simple roots, and, for $\alpha\in\Delta$,
\[
u_\alpha\in U_\alpha^\times(S),\qquad w_\alpha\in\uNorm_G(T)(S),
\]
the elements defined by the pinning.
